\documentclass[twocolumn]{aastex631}

\usepackage{amsmath}
\usepackage{multirow}
\usepackage{natbib}
\usepackage{graphicx} 
\usepackage{aas_macros}

\begin{document}

Horndeski theories, while constructed to be classically ghost-free, face the critical challenge of preserving this stability under quantum corrections, which often introduce problematic higher-derivative ghost terms into the effective action. This paper proposes the Quantum Kinetic Structure Preservation (QKSP) symmetry, a principle defined by specific algebraic and differential relations among the theory's functions, designed to explicitly maintain its intrinsic ghost-free kinetic structure at the one-loop level and prevent the generation of new propagating ghost degrees of freedom. Our multi-stage theoretical analysis first establishes classical stability conditions for Horndeski theories around a flat Friedmann-Lemaître-Robertson-Walker background, incorporating the observational constraint of luminal gravitational wave propagation ($c_T^2=1$). We then compute the one-loop effective action using the background field method and heat kernel regularization to identify quantum-generated higher-derivative terms. The QKSP conditions are derived by demanding the vanishing of coefficients for terms that would introduce higher-than-second-order time derivatives in the effective action for physical perturbations. We find that QKSP symmetry fundamentally implies that scalar and tensor modes must propagate at the same speed ($c_s^2 = c_T^2$), which, combined with the empirical $c_T^2=1$, leads to the strong prediction that QKSP-symmetric Horndeski theories must satisfy $c_s^2=1$. A subsequent phenomenological analysis of a representative QKSP-symmetric model, capable of mimicking the $\Lambda$CDM expansion history, reveals a gravitational slip parameter $\eta=1$ and time-dependent modifications to the effective gravitational constant, $G_{\text{eff}}$. For instance, a coupling parameter $|c_B|=0.1$ predicts percent-level deviations in $G_{\text{eff}}$ (e.g., a 7.4% enhancement or 6.4% suppression at $z=0$) that are within reach of current and future large-scale structure surveys. This work establishes a direct link between fundamental quantum stability requirements and observable cosmological phenomena, presenting a theoretically motivated and observationally falsifiable framework for modified gravity.

\end{document}
                